\chapter*{Preface}
\addcontentsline{toc}{chapter}{Preface}

Data science and machine learning are often organized around models. This text
organizes them around claims and the systems that carry those claims. Evidence
originates in files, instruments, databases, or services; software represents,
transforms, evaluates, and communicates it. A result is credible only to the
extent that its assumptions, implementation, provenance, uncertainty, and
operation can be examined.

Part I develops the computational foundations for that systems view. Part II
develops supervised learning from a precisely stated prediction contract through
linear models, classification, geometric methods, trees, ensembles, neural
networks, model selection, and a monitored release. The text
assumes mathematical maturity appropriate to graduate study, but not prior
professional experience with shells, virtual environments, packages, databases,
web APIs, continuous integration, or language-model agents. Tooling is taught
when it expresses a technical responsibility, not as an end in itself.

\section*{How this book relates to the repository}

The graduate text is the organizing artifact; the remaining repository
components make its arguments executable:

\begin{tabularx}{\textwidth}{@{}>{\bfseries}p{0.22\textwidth}X@{}}
\toprule
Textbook & Durable concepts, terminology, design arguments, connections, and
compact examples. \\
Lecture notebooks & Executable explanations, experiments, assertions, guided
practice, and extended reference material. \\
\code{rice\_dsm} package & Reusable code developed with professional structure,
documentation, types, exceptions, and tests. \\
Test suite and CI & Executable contracts that detect regressions across the
package, repository, and notebooks. \\
Supplementary guides & Slower introductions to tools and workflows that should
not be treated as assumed background. \\
\bottomrule
\end{tabularx}

Read the prose and notebooks together. Predict before execution, change one
assumption at a time, and state what the result does and does not establish.

\section*{A recurring distinction}

Throughout the course we separate three claims:

\begin{enumerate}
  \item \textbf{The code ran.} Python completed the requested operations.
  \item \textbf{The software contract holds.} Types, tests, schemas, and
  operational controls support the promised behavior.
  \item \textbf{The scientific conclusion is valid.} The data, assumptions,
  design, uncertainty, and interpretation justify the claim.
\end{enumerate}

None of these claims implies the next. A notebook can run while joining the
wrong records. A thoroughly tested implementation can faithfully compute an
invalid statistic. A sound analysis can still fail as a product if users cannot
reach it or operators cannot observe it.

\begin{professionalbox}
Professionalism is not ceremony added after discovery. It is the practice of
making assumptions, contracts, evidence, failure modes, and responsibility
visible early enough to improve the work.
\end{professionalbox}

\section*{Editorial status}

This is a developing graduate-course edition. Parts I and II now form connected
systems and supervised-learning arguments with worked examples, exercises,
appendices, and executable companions. Part II notebooks currently establish
the executable entry points and development contracts; their detailed
laboratories will be expanded unit by unit. The prose is ready for classroom
review while remaining open to correction, additional evidence, solutions, and
revision from use.
